\section{Three regimes for anchored gadgets}\label{sec:classification}

The closure law distinguishes sublinear from linear memory. In the gadget class, exact support and the anchors
also distinguish bounded one-round cost from cost that diverges. The mechanism is compactness at bounded
\emph{entropy}, followed by support-preserving conversion of a finite mixture to one maximally mixed bath.
Slofstra and Vidick read the same three-way split (bounded, finite but divergent, infinite) off hyperlinear
profiles~\cite[Section~2]{SlofstraVidick2018}; for sofic profiles of chunks, boundedness means that the chunk embeds in
a finite group~\cite[Fact~3.5]{Cornulier2013}. For a linear-system game that can be won with probability approaching one,
the entanglement needed to win with probability $1-\epsilon$ is governed by such a
profile~\cite[Corollaries~5.2 and~5.3]{SlofstraVidick2018}. It stays bounded exactly when a finite-dimensional perfect
strategy exists~\cite{CleveMittal2014,CleveLiuSlofstra2017}, and it is unbounded for some
games~\cite{Slofstra2019,Slofstra2020}. There the dimension diverges as the error tends to zero; here memory diverges
with the number of uses at fixed error, the leakage scale $1/n$ playing the role of the defect
(Remark~\ref{rem:hyperlinear-upper}). Channel profiles correspond to group profiles by~\cite[Theorem~9.2]{DouglasMghirbiB}.

\begin{theorem}[Gadget trichotomy]\label{thm:trichotomy}
For a canonical untwisted anchored gadget, exactly one of the following holds.
\begin{enumerate}[label=(\roman*),itemsep=3pt]
\item $\Phi$ has an exact finite tracial factorization. Equivalently, a finite-dimensional unitary
representation of $P_\Phi$ separates all labels, or a finite quotient
of $P_\Phi$ distinguishes every named label. Then $\sup_n L^*_\Phi(n)<\infty$ and
$\Qf(n,\epsilon)=O_\Phi(\log(n/\epsilon))$~\cite[Theorem~8.6]{DouglasMghirbiB}.
\item $\Phi\in\overline{\mathcal F}$ but has no exact finite tracial factorization.
Equivalently, the named labels are separated by a homomorphism
$P_\Phi\to U(\mathcal R^\omega)$ but by no finite quotient. Then
$L^*_\Phi(n)\to\infty$, $L^*_\Phi(n)=o(n)$, and $\Qs(n,\epsilon_n)\to\infty$ for every
$0<\epsilon_n\le1/16$. If also $\log(1/\epsilon_n)=o(n)$, flat-import rank-rate services have $Q=o(n)$.
\item $\Phi\notin\overline{\mathcal F}$. Equivalently, some distinct named
labels are identified by every homomorphism
$P_\Phi\to U(\mathcal R^\omega)$. Then $L^*_\Phi(n)=\Theta_\Phi(n)$ and
$\Qs(n,\epsilon_n)=\Theta_\Phi(n)$ for every error sequence $0<\epsilon_n\le1/16$.
\end{enumerate}
All upper constructions have $C=0$ and $P\le Q$.
\end{theorem}

In (i), \cite[Theorem~8.6]{DouglasMghirbiB} allows exchange $\lfloor n\log r/2\rfloor$. Skipping exports of at most
$\lceil\log r/2\rceil$ qubits, and holding from the start a maximally mixed qubit in place of each skipped import, meets the
budget $\lfloor(n-1)\log r/2\rfloor$ with the same error and $O_r(1)$ more memory.

Here $\mathcal R^\omega$ is a tracial ultrapower of the hyperfinite
$\mathrm{II}_1$ factor. Separation means operator-distinct images;
it does not require a map to preserve the regular group trace.
Padding by a trivial representation and tensor amplification turn
separation of the finite label set into the needed vanishing trace
overlaps (compare~\cite[Lemma~2.2 and the proof of Lemma~2.6]{SlofstraVidick2018}). Lemma~\ref{lem:ultrapower-labels} gives this conversion.
Finite-dimensional separation gives a finite quotient by residual
finiteness of the finitely generated matrix image. These criteria
concern the presented label group, not every element of its
source-group image.

In the middle regime memory diverges, possibly no faster than $\log n$, and the theorem gives no rate there.
In the linear regime the error threshold is the fixed value $1/16$, whereas the general converse of
Theorem~\ref{thm:closure-law} needs error below the channel's distance from the closure.

\paragraph{Why bounded entropy is enough.}
Take rounds with bounded source entropy and production $a_j\to0$. Bounded entropy keeps the sorted
source spectra from spreading over ever more eigenvalues, so they converge. A round moves weight
between bath eigenvalues through a doubly stochastic matrix, and small production forces the spectrum
after a use to match the spectrum before it, up to $\sqrt{2\ln2\,a_j}$ in $\ell^1$. In the limit the
round maps each eigenspace of the limiting source spectrum unitarily onto itself. So the limiting channel
is a finite mixture of finite tracial factorizations (Lemma~\ref{lem:entropy-compactness}).
This mixture conclusion holds for general channels. Anchors turn an exact-support mixture close to
the canonical gadget into a separating representation.

Let $B=\sup_n L^*_\Phi(n)$. Let $R_{\rm sep}$ be the least $\log m$ for which a representation $\pi$ of
$P_\Phi$ on $\C^m$ satisfies $|\tau_m(\pi(v_g)^*\pi(v_h))|\le1/2$ for all distinct labels $g\ne h$. Then
\[
 R_{\rm sep}\le16(B+1)+\log_2(33(d^4+1)),\qquad
 B\le\lceil\log_2(8r(r-1))\rceil R_{\rm sep}.
\]
The exact finite tracial factorization is obtained by Malcev residual finiteness of the finitely generated
matrix image and its finite quotient's regular representation~\cite{Nica2013}. No bound on that quotient's
order follows from the displayed representation bound. Appendix~\ref{app:compactness} proves the classification.

\subsection{Known examples}

\begin{table}[ht]
\centering\small
\begin{tabularx}{\textwidth}{@{}>{\raggedright\arraybackslash}p{0.2\textwidth}>{\raggedright\arraybackslash}X>{\raggedright\arraybackslash}X@{}}
\toprule
Regime & Criterion on $P_\Phi$; memory & Known examples\\
\midrule
(i) exact finite tracial factorization
 & A finite quotient separates the labels; memory $O(\log(n/\epsilon))$.
 & Labels separated in a finite, residually finite or maximally almost periodic group satisfying the
   relations (Corollary~\ref{cor:map}): the two-element example of Section~\ref{sec:gadgets}, the
   lamplighter $C_2\wr\Z$.\\
\addlinespace
(ii) in $\overline{\mathcal F}$, no finite factorization
 & Labels separated in $U(\mathcal R^\omega)$ but by no finite quotient; memory diverges and is $o(n)$.
 & Houghton's group $H_3$; the configuration-lamp group $\Gamma$.\\
\addlinespace
(iii) outside $\overline{\mathcal F}$
 & Two labels identified by every homomorphism into $U(\mathcal R^\omega)$; memory $\Theta(n)$.
 & Known only conditionally: the canonical gadget $\Phi_{\rm cert}$ of the group certificate
   of~\cite{MghirbiDouglasD}, at an explicit distance from $\overline{\mathcal F}$
   (Proposition~\ref{prop:explicit-canonical}). Nonempty exactly when some finitely presented group
   is not hyperlinear (Proposition~\ref{prop:regime-iii}).\\
\bottomrule
\end{tabularx}
\caption{The three regimes of Theorem~\ref{thm:trichotomy}, at fixed error at most $1/16$.}\label{tab:regimes}
\end{table}

The placements in regimes (i) and (ii) follow from results in this paper; regime (iii) is
discussed after Proposition~\ref{prop:regime-iii}. The two-element example
has $P_\Phi=\Z/2$. A gadget from the lamplighter $C_2\wr\Z$ has labels separated by finite quotients
(Section~\ref{sec:lower}). The Houghton gadget has $P_{\Phi_H}\cong H_3$ (Section~\ref{sec:gadgets}).
Every finite-dimensional representation of $H_3$ identifies a three-cycle label with the label $1$
(Remark~\ref{rem:houghton-map}), so regime (i) fails. Since $H_3$ is amenable, it is hyperlinear, so it
embeds in $U(\mathcal R^\omega)$~\cite{Pestov2008} and the labels are separated: this is regime (ii).
For the configuration-lamp gadget, Corollary~\ref{cor:lamp} gives $L^*_\Phi(n)\to\infty$ and
$\Qs(n,\epsilon)=o(n)$ at fixed error, which excludes regimes (i) and (iii).

\begin{proposition}[Regime (iii) and non-hyperlinear groups]\label{prop:regime-iii}
The following are equivalent.
\begin{enumerate}[label=(\alph*),itemsep=1pt]
\item Some canonical untwisted anchored gadget lies in regime (iii) of Theorem~\ref{thm:trichotomy}.
\item Some finitely presented group $P$ has an element $\zeta\ne1$ such that $\pi(\zeta)=1$ for every
homomorphism $\pi:P\to U(\mathcal R^\omega)$.
\item Some finitely presented group is not hyperlinear.
\end{enumerate}
\end{proposition}

For (a)$\Rightarrow$(b), take $P=P_\Phi$ and $\zeta=v_g^{-1}v_h$ for two identified labels. For
(b)$\Rightarrow$(a), add the freely trivial word $ww^{-1}$, where $w$ represents $\zeta$, to a
presentation of $P$; this makes $\zeta$ a label with its own anchor. Appendix~\ref{app:regime-iii}
gives the proof, including the standard step (c)$\Rightarrow$(b).

\paragraph{Which gadgets are known outside the closure?}
Our companion paper~\cite{MghirbiDouglasD} gives explicit rational channels outside
$\overline{\mathcal F}$, at explicit distances. Each factorizes exactly through a group von Neumann
algebra. Its distances are the normalized Choi and half-diamond distances of Section~\ref{sec:model},
and its closure of the noisy operations is $\overline{\mathcal F}$, so its bounds apply here unchanged.
Write $\operatorname{dist}_\diamond(\Phi,\overline{\mathcal F})=\inf_{\Psi\in\overline{\mathcal F}}\ddia(\Phi,\Psi)$.

\begin{proposition}[Mghirbi--Douglas~{\cite[Theorems~I and~III]{MghirbiDouglasD}}]\label{prop:rational-distances}
Put $\delta_{\rm WZ}=2^{-2^{2^{1122000}}}$, the error tolerance of Wang and
Zhi~\cite[(2.20)]{WangZhi2026}, and assume the theorems of~\cite{WangZhi2026}.
\begin{enumerate}[label=(\alph*),itemsep=1pt]
\item A unital channel $\Phi_{12}$ on twelve qubits, with Kraus entries in $\{0,1,\pm3/5,4/5\}$ and
Choi rank $949$, has $\operatorname{dist}_J(\Phi_{12},\overline{\mathcal F})\ge2^{-50}\delta_{\rm WZ}^2$
and $\operatorname{dist}_\diamond(\Phi_{12},\overline{\mathcal F})\ge2^{-50}\delta_{\rm WZ}^2$.
\item The fifteen-qubit channel $\Phi_{15}$ and its reductions $\Phi_1,\Phi_2$, on $M_d$ with
$d=31\,981$ and $31\,982$, have Choi rank $11\,230$ and both distances at least
$2^{-41}\delta_{\rm WZ}^2$.
\item The fully referenced channel $\Phi_{\rm full}$ on $M_{44\,086}$, with a singleton anchor for
each of its $11\,230$ labels, has
$\operatorname{dist}_J(\Phi_{\rm full},\overline{\mathcal F})\ge2^{-36}\delta_{\rm WZ}^2$ and
$\operatorname{dist}_\diamond(\Phi_{\rm full},\overline{\mathcal F})\ge2^{-21}\delta_{\rm WZ}^2$.
\end{enumerate}
\end{proposition}

The constants are astronomically small: already $2^{1122000}$ has $337\,756$ decimal digits. They
come from the rigidity estimate, not from the channels. Theorem~\ref{thm:closure-law}(b) turns them
into explicit purity and memory bounds, as~\cite[Section~6]{MghirbiDouglasD} notes.

\begin{corollary}[Theorem~\ref{thm:closure-law}(b) with~{\cite[Theorem~III]{MghirbiDouglasD}}]\label{cor:rational-memory}
Every between-round service of $\Phi_{\rm full}$ with error at most $\epsilon$ satisfies
\[
 P+C\ge\frac{n\bigl(2^{-36}\delta_{\rm WZ}^2-\epsilon\bigr)_+^2}{2\ln2},\qquad
 Q+C\ge\frac{n\bigl(2^{-36}\delta_{\rm WZ}^2-\epsilon\bigr)_+^2}{2\ln2}.
\]
With maximally mixed imports and $\epsilon\le2^{-37}\delta_{\rm WZ}^2$, this gives
$Q\ge2^{-75}\delta_{\rm WZ}^4\,n$.
\end{corollary}

The memory is linear, but the bound reaches one qubit only at $n=2^{75}\delta_{\rm WZ}^{-4}$. It
replaces a compactness argument by a number; it does not describe an experiment.

None of these channels is a canonical gadget, so Theorem~\ref{thm:trichotomy} does not apply to
them. Their blocks use $3/5$ and $4/5$ where~\eqref{eq:canonical-unitary} uses $1/\sqrt2$. Equal
triangles share one block, while the canonical table gives every tagged prefix its own. Most labels
have no anchor. $\Phi_{12}$, the output of the anchored-letter
compiler~\cite[Definition~3.1]{MghirbiDouglasD}, has singletons only for the identity, the letters and
some further labels: $364$ of its $949$ labels. $\Phi_2$ and $\Phi_{15}$ have singletons only for the
identity and the witness, and $\Phi_1$ only for the identity. $\Phi_{\rm full}$ anchors every label but
keeps the $3/5,4/5$ blocks.

A canonical example comes from the same group construction. Take the presentation
of~\cite[Section~7]{MghirbiDouglasD}, with $43$ generators and $4\,417$ relators of length at most
$19$, and its group image $\Gamma$. One generator is the witness $j$, and $j\ne1$ in $\Gamma$. Let
$\Phi_{\rm cert}$ be the canonical gadget of $\Gamma$, these generators and these relators. It acts on
$M_d$ with $d=108\,364$, and the letter $j$ has its own singleton anchor. Every homomorphism from the
presented group into a tracial matrix ultraproduct kills $j$~\cite[Theorem~1]{MghirbiDouglasDv1}. The
presented group is countable, so a label-separating homomorphism into $U(\mathcal R^\omega)$ would give
one into a tracial matrix ultraproduct: the von Neumann algebra generated by its image embeds
trace-preservingly into one. That would contradict the certificate, so
Lemma~\ref{lem:ultrapower-labels} places $\Phi_{\rm cert}$ outside $\overline{\mathcal F}$, in
regime~(iii).

The fully referenced bound of~\cite{MghirbiDouglasD} makes this example explicit.

\begin{proposition}[An explicit canonical gadget outside the closure]\label{prop:explicit-canonical}
Assume the rigidity theorem~\cite[Theorem~D.4]{MghirbiDouglasD}, which rests on~\cite{WangZhi2026}.
Then
\begin{align*}
 \operatorname{dist}_J(\Phi_{\rm cert},\overline{\mathcal F})&\ge\frac{\delta_{\rm WZ}^2}{4096\cdot19^2\cdot108\,364}>2^{-38}\delta_{\rm WZ}^2,\\
 \operatorname{dist}_\diamond(\Phi_{\rm cert},\overline{\mathcal F})&\ge\frac{\delta_{\rm WZ}^2}{4096\cdot19^2}>2^{-21}\delta_{\rm WZ}^2.
\end{align*}
\end{proposition}

\begin{proof}
By~\cite[Theorem~D.4]{MghirbiDouglasD}, unitaries whose relator defects are at most
$\delta_{\rm WZ}$ have $\|v_j-I\|_2<2^{-17}$, in every matrix dimension. So $\delta_{\rm WZ}$ is a
relator tolerance in the sense of~\cite[Lemma~7.3]{MghirbiDouglasD}. That lemma, from~\cite[Theorem~6]{MghirbiDouglasDv1},
gives the displayed bounds, with $\ell_*=19$, for channels built from $3/5,4/5$ blocks with a
singleton for every label. Its proof, with $k$ as there, uses the coefficients only twice. Each block
entry $c$ has $|c|\le4/5$, which keeps the bath block within $3k$ of $c$ times its label unitary. Each
block has $ab\ge12/25$, which turns the row-orthogonality defect, $13k$ or $(12+\sqrt2)k$, into a
triangle defect below $32k$. The canonical coefficients $a=b=1/\sqrt2$ satisfy both. The canonical
table contains every prefix triangle and every inverse triangle, every label has an anchor, and the
anchors of $1$ and $j$ are distinct. The proof uses one anchor per label and treats blocks one at a
time, so tagged repeats change nothing.
\end{proof}

So $\Phi_{\rm cert}$ is a canonical gadget in regime~(iii) at an explicit distance.
Theorem~\ref{thm:closure-law}(b) gives
$P+C\ge n(2^{-38}\delta_{\rm WZ}^2-\epsilon)_+^2/(2\ln2)$ for its between-round services, and
Theorem~\ref{thm:trichotomy}(iii) gives linear memory at every error up to $1/16$, with a constant
from compactness.

These results are conditional. The certificate rests on Alekseev--Liu--Thom, Theorem~B, and Thom,
Theorem~1.2~\cite{AlekseevLiuThom2026,Thom2026}. The explicit distances rest on the rigidity theorems
of Wang and Zhi~\cite{WangZhi2026}, whose proofs follow steps of those two
papers;~\cite[Section~1.4]{MghirbiDouglasD} lists which statements use them. All three are recent
unrefereed preprints at the time of writing. Exact factorizability is unconditional. By contrast
$\mathrm{MIP}^*=\mathrm{RE}$ and Haagerup--Musat imply nonconstructively an exactly factorizable
nonclosure channel in some dimension, recorded for a Schur channel
in~\cite[Corollary~3.11]{Musat2025}, but do not by themselves supply an explicit group or canonical
gadget.
